\section{Verification Results}
\label{sec:results}

Table~\ref{tab:results} is the complete result matrix, rendered from
\texttt{results/results.json}. Every cell is populated: proved, refuted,
inconclusive, or non-terminating with a recorded cause and the budget it
consumed. Of \wgProps{} properties, \wgProved{} are proved, \wgAttacks{} are
refuted, \wgInconcl{} is inconclusive and \wgNonterm{} did not terminate.

\begin{table*}[t]
\centering
\caption{Verification outcomes, the complete matrix. \wgProVerif{} and
\wgTamarin, single core, \wgMemCap\,GB address-space cap, \wgBudget\,s
wall-clock budget. \(\dagger\) marks a budget consumed without reaching a
verdict, which is not a time-to-verdict and must not be read as one; \emph{n/r}
marks a run that predates per-query cost recording. Every cell is rendered
from \texttt{results/results.json}.}
\label{tab:results}
\footnotesize
\begin{tabular}{@{}lllrl@{}}
\toprule
ID & Property & Prover & Cost & Outcome \\
\midrule
P1 & Session-key secrecy & ProVerif & \textit{n/r} & \ok \\
P3 & Forward secrecy & ProVerif & \textit{n/r} & \ok \\
P4 & Post-quantum forward secrecy (psk) & ProVerif & \textit{n/r} & \ok \\
P4c & \quad same, psk absent (control) & ProVerif & \textit{n/r} & \attack \\
P5a & Agreement, responder to initiator & ProVerif & \textit{n/r} & \nonterm \\
P5b & Agreement, initiator to responder & ProVerif & \textit{n/r} & \ok \\
P5c & Injective agreement / replay (P9) & ProVerif & \textit{n/r} & \ok \\
P6a & KCI, commitment level & ProVerif & \textit{n/r} & \nonterm \\
P6b & KCI, completion level & ProVerif & \textit{n/r} & \nonterm \\
P12a & KEM-derived psk secrecy, K-PK binding & ProVerif & \textit{n/r} & \ok \\
P12b & \quad same, binding absent (control) & ProVerif & \textit{n/r} & \attack \\
P12c & KEM-derived psk, responder-view soundness & ProVerif & \textit{n/r} & \inconcl \\
\midrule
S0 & Model admits an honest run & Tamarin & \textit{n/r} & \ok \\
S1 & Session-key secrecy & Tamarin & \textit{n/r} & \ok \\
S3 & Forward secrecy & Tamarin & \textit{n/r} & \ok \\
S5 & Agreement, no compromise & Tamarin & \textit{n/r} & \ok \\
S6a & KCI, commitment level & Tamarin & \textit{n/r} & \attack \\
S6b & KCI, completion level & Tamarin & \textit{n/r} & \ok \\
\bottomrule
\end{tabular}

\end{table*}

\subsection{Baseline secrecy and authentication}

\para{Secrecy.} Session-key secrecy (P1) and forward secrecy (P3) are proved
with an unbounded number of concurrent sessions. Forward secrecy is stated by
releasing both static scalars in phase~1, strictly after every phase-0
session; the ephemeral-ephemeral term \(ee\) is what survives.

\para{Authentication, level by level.} Non-injective agreement of the
responder with the initiator (P5b) is proved, and so is injective agreement
(P5c). Injectivity is simultaneously the replay-resistance statement: a
replayed initiation would produce two responder completions against a single
initiator run, which is exactly what the injective query forbids, and the
TAI64N timestamp is what prevents it.

\para{The direction that resists.} Agreement in the opposite direction (P5a)
did not terminate. Three formulations were attempted---agreement on both
static values and both transport keys, an existentially weakened aliveness
statement, and the same with \texttt{redundantHypElim} enabled---and each
exhausted the budget. The direction that resists is the one whose proof must
reason through the authenticated decryption of message~2 under a psk-derived
key, the deepest point in the key schedule.

\subsection{What the pre-shared key actually contributes}

\begin{table}[t]
\centering
\caption{The post-quantum control experiment: two models, one variable. Both
arms face the same attacker, which gains a discrete-logarithm oracle in
phase~1. The verdict column is read from the same result file as
Table~\ref{tab:results}, so the experiment cannot be reported as having
succeeded in a run where it did not.}
\label{tab:psk}
\footnotesize
\begin{tabular}{@{}llc@{}}
\toprule
Model & psk & Transport-key secrecy \\
\midrule
\texttt{22\_pq\_forward\_secrecy.pv} & secret & \ok \\
\texttt{23\_pq\_fs\_nopsk.pv} & public & \attack \\
\bottomrule
\end{tabular}

\end{table}

The claim that \texttt{IKpsk2}'s pre-shared key provides post-quantum
protection is common and rarely tested. Table~\ref{tab:psk} tests it. Both
models are the same file up to one line: whether the pre-shared key is
published. Both face an attacker that, in phase~1, can invert every
Curve25519 public value it recorded---which is every ephemeral, since they
travel in clear, and both statics.

\para{With the key public.} Every input to the key schedule is recoverable and
ProVerif returns a derivation reconstructing both transport keys; the
derivation is archived under \texttt{results/traces/}.

\para{With the key secret.} The same attacker is defeated:
\(\mathrm{KDF}_3\) mixes an independent secret into the chaining key, and no
amount of discrete-logarithm capability recovers it. Because the two models
differ in exactly one variable, the difference in outcome is attributable to
that variable and to nothing else. This is the sense in which the pre-shared
key is doing post-quantum work, and the control is what makes the statement
more than an assertion.

\subsection{Where KCI resistance resides}

\begin{figure}[t]
\centering
\resizebox{\columnwidth}{!}{\begin{tikzpicture}[x=1cm,y=1cm,font=\footnotesize]
  \node[draw=badred,thick,rounded corners,align=left,inner sep=6pt,
        text width=7.0cm] (a) at (0,0)
    {\textbf{\color{badred}1. Forge message 1}\\[2pt]
     needs $es=\mathrm{DH}(e_i,S_r)$ and $ss=\mathrm{DH}(s_i,S_r)$\\[1pt]
     {\color{black!55}both computable from the leaked responder scalar}\\[3pt]
     \hrule\vspace{3pt}
     {\color{badred}\textbf{S6a falsified}} --- the responder commits key
     material to a session it attributes to an honest initiator};
  \node[draw=okgreen,thick,rounded corners,align=left,inner sep=6pt,
        text width=7.0cm] (b) at (0,-3.15)
    {\textbf{\color{okgreen}2. Complete the session}\\[2pt]
     additionally needs $se=\mathrm{DH}(e_r,S_i)$\\[1pt]
     {\color{black!55}requires the initiator's static scalar, which is not
      compromised}\\[3pt]
     \hrule\vspace{3pt}
     {\color{okgreen}\textbf{S6b verified}} --- no transport record is ever
     accepted on that session};
  \draw[->,very thick,black!45] (a.south) -- node[right,black!55,align=left]
    {\scriptsize the gap $se$ holds open} (b.north);
\end{tikzpicture}
}
\caption{With the responder's static key compromised, two adjacent
authentication levels separate. The static-ephemeral term \(se\) is the whole
of the difference, and the property lives in the gap between them.}
\label{fig:kci}
\end{figure}

Key-compromise impersonation asks whether compromising a party's own long-term
key lets an attacker impersonate \emph{others} to it. We asked it at two
levels, and the levels separate (Figure~\ref{fig:kci}).

\para{The responder does commit.} Producing a well-formed message~1 in the
name of an honest initiator requires \(es=\mathrm{DH}(e_i,S_r)\) and
\(ss=\mathrm{DH}(s_i,S_r)\). Both are computable from the leaked responder
scalar, so an attacker can drive the responder to commit transport keys for a
session it attributes to that initiator. Tamarin refutes the
commitment-level lemma (S6a) and returns the trace.

\para{The responder never completes.} Deriving those transport keys
additionally requires \(se=\mathrm{DH}(e_r,S_i)\), and that needs either the
responder's ephemeral scalar or the initiator's static scalar. The attacker
has neither. Tamarin verifies the completion-level lemma (S6b): the responder
never accepts a transport record on such a session.

\para{Why the distinction is operational.} WireGuard's KCI resistance depends
on treating a session as established only once authenticated traffic has
arrived, not once the handshake response has been sent. A model with a single
``session established'' event would report one verdict where there are two,
and would locate the property in the wrong place.

\subsection{KEM-derived pre-shared keys}

If the pre-shared key is to come from a KEM rather than from manual
configuration---the construction Rosenpass~\cite{rosenpass2023} uses---the
question is what that KEM must bind. We modelled two: one whose shared secret
depends on the encapsulation key, and one whose shared secret depends only on
the encapsulation randomness, together with the re-encapsulation capability
that absence of binding confers.

\para{Binding is the whole of it.} With binding, secrecy of the established
pre-shared key is proved (P12a). Without it, the attacker redirects an
observed ciphertext to an encapsulation key of its own and recovers the
pre-shared key outright (P12b). The failure is not a weakness in any
primitive: confidentiality and integrity of the KEM are untouched, and only
the binding is missing. In the vocabulary of Cremers, Dax and
Medinger~\cite{cremers2024kembinding} the separating property is
MAL-BIND-K-PK, which ML-KEM satisfies.

\para{One query we could not settle.} A third query, asking whether the
responder's view of the established key is sound, returned inconclusive
(P12c): a bare encapsulation authenticates nobody, and the honest statement is
disjunctive in a way ProVerif's approximation could not decide. We report it
as inconclusive rather than restate it until it passed. An undecided query is
a result; a restated one is not.
